\subsubsection{同一代理下的优化不足与目标不足}
原库与重拟合 MILP 使用同一个冻结代理，不能从主问题设计域更大推出已取得动作更好。
独立重算 $255$ 条双方完成的历史：库动作的 QP 下界高于 MILP 所选动作上界 $99$ 条，反向 $16$ 条，区间重叠 $140$ 条；另 $1$ 条库选址失败保留在总分母中。
全部 $515$ 次原平方弦上界主问题均超时，未得到 incumbent 或 global best-bound，故未收紧初始全域上界。
$256$ 次选址中，$128$ 次满足冻结的总代价相对 gap，另外 $128$ 次停于预算或迭代限制。
前一组仍有 $43$ 条相对恒充的动作增益区间包含零；历史残差底数较大时，相对总代理代价停止不能自动分辨新增收益。
双方唯一正确为 $67$ 条，库单独 $14$ 条、MILP 单独 $4$ 条。
库单独正确的 $14$ 条中，$6$ 条代理区间严格库优，$8$ 条重叠；MILP 单独正确的 $4$ 条中，反有 $1$ 条严格库代理优。
因此现有差异同时显示计算限制和代理排序不足；尚不能把净差 $10$ 条逐一归因于线性化、解释域或门控。
完整逐历史区间与停止记录见 \path{evidence/adaptive_C_optimization_diagnostic/}。
